Um zu zeigen, dass die Wahrscheinlichkeitsdichte \(\rho = \Psi^* \Psi\) die Kontinuitätsgleichung erfüllt, beginnen wir mit der Schrödinger-Gleichung und ihrer komplex konjugierten Form:

\[
\left\{-\frac{\hbar^2 \nabla^2}{2 m} + V(\mathbf{x})\right\} \Psi = i \hbar \frac{\partial \Psi}{\partial t}
\]

\[
\left\{-\frac{\hbar^2 \nabla^2}{2 m} + V(\mathbf{x})\right\} \Psi^* = -i \hbar \frac{\partial \Psi^*}{\partial t}
\]

Multiplizieren wir die Schrödinger-Gleichung mit \(\Psi^*\) und die konjugierte Schrödinger-Gleichung mit \(\Psi\):

\[
\Psi^* \left\{-\frac{\hbar^2 \nabla^2}{2 m} \Psi + V(\mathbf{x}) \Psi\right\} = \Psi^* i \hbar \frac{\partial \Psi}{\partial t}
\]

\[
\Psi \left\{-\frac{\hbar^2 \nabla^2}{2 m} \Psi^* + V(\mathbf{x}) \Psi^*\right\} = -\Psi i \hbar \frac{\partial \Psi^*}{\partial t}
\]

Subtrahieren wir die zweite Gleichung von der ersten:

\[
\Psi^* \left(-\frac{\hbar^2 \nabla^2}{2 m} \Psi\right) - \Psi \left(-\frac{\hbar^2 \nabla^2}{2 m} \Psi^*\right) = \Psi^* i \hbar \frac{\partial \Psi}{\partial t} + \Psi i \hbar \frac{\partial \Psi^*}{\partial t}
\]

Vereinfachen wir die linke Seite:

\[
-\frac{\hbar^2}{2 m} \left(\Psi^* \nabla^2 \Psi - \Psi \nabla^2 \Psi^*\right) = i \hbar \left(\Psi^* \frac{\partial \Psi}{\partial t} + \Psi \frac{\partial \Psi^*}{\partial t}\right)
\]

Die linke Seite kann umgeschrieben werden als:

\[
\nabla \cdot \left(\frac{\hbar}{2 i m} \left(\Psi^* \nabla \Psi - \Psi \nabla \Psi^*\right)\right)
\]

Diese Umformung zur Divergenzform ist korrekt, da der Ausdruck \(\Psi^* \nabla^2 \Psi - \Psi \nabla^2 \Psi^*\) durch Anwendung der Produktregel umgeschrieben werden kann. Die Produktregel besagt, dass \(\nabla \cdot (\Psi^* \nabla \Psi) = \Psi^* \nabla^2 \Psi + \nabla \Psi^* \cdot \nabla \Psi\) und \(\nabla \cdot (\Psi \nabla \Psi^*) = \Psi \nabla^2 \Psi^* + \nabla \Psi \cdot \nabla \Psi^*\). Subtrahiert man diese, erhält man die Divergenzform \(\nabla \cdot (\Psi^* \nabla \Psi - \Psi \nabla \Psi^*)\).

Setzen wir dies in die Kontinuitätsgleichung ein:

\[
\frac{\partial}{\partial t} (\Psi^* \Psi) + \nabla \cdot \left(\frac{\hbar}{2 i m} \left(\Psi^* \nabla \Psi - \Psi \nabla \Psi^*\right)\right) = 0
\]

Dies zeigt, dass die Wahrscheinlichkeitsdichte \(\rho = \Psi^* \Psi\) die Kontinuitätsgleichung erfüllt, wobei der Stromdichtevektor \(\mathbf{j}\) gegeben ist durch:

\[
\mathbf{j} = \frac{\hbar}{2 i m} \left(\Psi^* \nabla \Psi - \Psi \nabla \Psi^*\right)
\]